%*******************************************************
% Conclusion generale
% Chapitre non numerote, place apres le chapitre Perspectives.
%*******************************************************
\pdfbookmark[1]{Conclusion générale}{ConclusionGenerale}
\phantomsection
\addcontentsline{toc}{chapter}{\tocEntry{Conclusion générale}}

\chapter*{Conclusion générale}
\label{ch:conclusion}
\markboth{Conclusion générale}{Conclusion générale}

% ====================================================================================
% Longueur visee: 3 pages. Celle de Chatelain en fait 2.
%
% A ECRIRE EN MEME TEMPS QUE \cref{ch:introduction}, d'une seule traite: le role
% de la conclusion est de refermer ce que l'introduction a ouvert. Ecrites a
% quelques semaines d'intervalle, elles se repetent au lieu de se repondre.
%
% CE QUE CE CHAPITRE N'EST PAS: un resume chapitre par chapitre. Les bilans de
% \cref{ch:robust-detection} et \cref{ch:spdnet} ont deja fait ce travail, et le
% refaire ici est la facon la plus sure de perdre le lecteur a la derniere page.
% ====================================================================================

% ============ TODO texte (~1 page) --- CE QUE LES CHAPITRES ETABLISSENT ENSEMBLE ===
% Ce que les trois chapitres de contribution etablissent CONJOINTEMENT, et
% qu'aucun n'etablit seul.
%
% Piste: dans les trois cas, la contrainte --- statistique au
% \cref{ch:robust-detection}, geometrique au \cref{ch:spdnet} --- n'est pas un
% raffinement de performance mais ce qui rend le probleme faisable au nombre
% d'observations reellement disponible. C'est la meme phrase que le bilan du
% \cref{ch:robust-detection}, mais elevee d'un cran: elle vaut aussi quand la
% contrainte est apprise et non imposee.
% ====================================================================================

% ============ TODO texte (~0,5 page) --- LA THESE QUE L'ON DEFENDRAIT ==============
% L'unique affirmation methodologique que l'on defendrait a la soutenance:
% choisir une moyenne, c'est choisir une geometrie, donc un modele statistique.
%
% La reformuler autrement qu'en \cref{ch:introduction} et qu'en
% \cref{prop:spdnet-moyennes-frechet} --- ici, ce n'est plus une annonce ni une
% proposition, c'est un bilan: on a vu ce que cela coute et ce que cela achete.
% ====================================================================================

% ============ TODO texte (~0,5 page) --- LA REPRODUCTIBILITE COMME RESULTAT ========
% Ce que l'appareil de reproductibilite demontre en tant que METHODE DE TRAVAIL,
% et non comme detail technique.
%
% L'argument le plus fort est factuel: c'est en instrumentant les simulations
% que des erreurs de signe et de conjugaison presentes dans du code deja publie
% sont apparues. Le dire sobrement, sans en faire une lecon.
% ====================================================================================

% ============ TODO texte (~0,5 page) --- CE QUE LE MEMOIRE NE TRANCHE PAS ==========
% Un jury d'hdr attend cette section, et elle se lit comme de la maturite et non
% comme un aveu. Points disponibles, tous deja etablis dans le corps du texte:
%
%  - l'unicite de l'estimateur a deux textures n'est etablie que numeriquement
%    (\cref{eq:detection-mills-tyler}); la convergence sous-lineaire mesuree est
%    un RESULTAT, pas un dysfonctionnement --- le dire ainsi;
%  - les seuils en forme close manquent pour les detecteurs structures
%    adaptatifs, d'ou le recours systematique a la simulation;
%  - aucun des travaux ne combine estimation en ligne ET contrainte de structure
%    (Toeplitz), la contrainte n'etant pas compatible avec la geometrie du cone;
%  - une anomalie non expliquee subsiste sur la puissance en regime gaussien
%    (non-monotonie a T faible). Si elle est mentionnee en legende, y renvoyer;
%    sinon, ne pas l'introduire ici.
% ====================================================================================

% ============ TODO texte (~4 lignes) --- PASSAGE DE RELAIS ========================
% Une transition courte vers \cref{ch:perspectives}, de sorte que les deux
% chapitres ne se recouvrent pas: la conclusion dit ce qui est acquis et ce qui
% ne l'est pas, les perspectives disent ce que l'on va faire. Si les deux
% chapitres disent la meme chose, c'est celui-ci qu'il faut raccourcir.
% ====================================================================================
